\documentclass[../../../hsm_tok_q.tex]{subfiles}

\begin{document}
    \setcounter{figure}{0}
    \textsf{図\ref{fig/hsm_tok_2025_1st_05_01_q}}に示した立体ABCD--EFGHは，\\
    $\mathrm{AB}=\mathrm{AD}=6\,\mathrm{cm}$，$\mathrm{AE}=4\,\mathrm{cm}$の直方体である．

    頂点Aと頂点Cを結び，線分AC上にある点をPとする．

    頂点Hと点Pを結ぶ．

    次の各問に答えよ．
    %
    \vspace{.5\baselineskip}
    {\par \centering
        \setlength{\abovecaptionskip}{3pt}
        \includegraphics%
            {\subfix{fig_hsm_tok_2025_1st_05/fig_hsm_tok_2025_1st_05_01_q.pdf}}
            \captionof{figure}{} \label{fig/hsm_tok_2025_1st_05_01_q}
    \par}
    \vspace{.4\baselineskip}
    %
    \begin{enumerate}
        \item 次の\:\myboxf{　}\:の中の「\textsf{お}」「\textsf{か}」に当てはまる%
            数字をそれぞれ答えよ．

            \textsf{図\ref{fig/hsm_tok_2025_1st_05_01_q}}において，%
            頂点Dと点P，頂点Eと点Pをそれぞれ結んだ場合を考える．

            点Pが線分ACの中点のとき，立体P--AEHDの体積は，%
            \:\myboxf{おか}\:$\mathrm{cm}^3$である．
        \newpage
        \item 次の\:\myboxf{　}\:の中の「\textsf{き}」「\textsf{く}」「\textsf{け}」に当てはまる
            数字をそれぞれ答えよ．

            \textsf{図\ref{fig/hsm_tok_2025_1st_05_02_q}}は，%
            \textsf{図\ref{fig/hsm_tok_2025_1st_05_01_q}}において，%
            頂点Fと頂点H，頂点Fと点Pをそれぞれ結んだ場合を表している．

            $\mathrm{AP}:\mathrm{PC}=5:1$ のとき，\\           
            $\triangle\mathrm{FPH}$の面積は，%
            $\myboxf{きく}\,\sqrt{\,\myboxf{け}\,}\:\mathrm{cm}^2$である．
            \\
            \begin{minipage}{\linewidth}
                \vspace{.5\baselineskip}
                {\par \centering
                    \setlength{\abovecaptionskip}{3pt}
                    \includegraphics%
                        {\subfix{fig_hsm_tok_2025_1st_05/fig_hsm_tok_2025_1st_05_02_q.pdf}}
                        \captionof{figure}{} \label{fig/hsm_tok_2025_1st_05_02_q}
                \par}
            \end{minipage}
    \end{enumerate}
\end{document}
